% =====================================================================
% 05-Fusion.tex
% Geometric Hallucination Taxonomy — Mathematical Spine
% Borealis S. Hedling — 2026
% =====================================================================

\section*{05. Fusion–Field Operator}

This module formalizes nonlinear fusion fields, interaction terms,
dominance conditions, and instability criteria used to define
Fusion–Field Hallucinations (FFH) in the Geometric Hallucination
Taxonomy (GHT). All definitions rely on the geometric foundations
established in 00-Foundations.tex.

Fusion–field instability corresponds to uncontrolled nonlinear blending
of latent representations, producing representational spikes, semantic
cross–contamination, and collapse into mixed manifolds.

% ---------------------------------------------------------------------
\subsection*{5.1 Fusion–Field Definition}

Let $x, y \in T_p M$ be tangent vectors representing latent states.
A fusion–field is a smooth map
\[
F : TM \times TM \to TM
\]
defined by
\[
F(x,y)
= \alpha x + \beta y + \gamma \, \sigma(x,y),
\]
where:
\begin{itemize}
\item $\alpha, \beta, \gamma \in \mathbb{R}$ are scalar weights,
\item $\sigma : TM \times TM \to TM$ is a nonlinear interaction term.
\end{itemize}

The linear terms $\alpha x$ and $\beta y$ represent stable blending.
The nonlinear term $\gamma \, \sigma(x,y)$ represents interaction
dynamics that may induce instability.

% ---------------------------------------------------------------------
\subsection*{5.2 Properties of the Interaction Term}

The interaction term $\sigma$ satisfies:
\begin{itemize}
\item smoothness: $\sigma$ is $C^\infty$,
\item bilinearity in the linear regime,
\item nonlinear amplification in unstable regimes,
\item symmetry or antisymmetry depending on the operator subtype.
\end{itemize}

Examples include:
\[
\sigma(x,y) = \langle x,y \rangle_g \, x,
\qquad
\sigma(x,y) = x \times y,
\qquad
\sigma(x,y) = \mathrm{ReLU}(x+y),
\]
depending on the latent geometry.

% ---------------------------------------------------------------------
\subsection*{5.3 Fusion–Field Norm}

The norm of the fusion–field output is
\[
\|F(x,y)\|_g
= \|\alpha x + \beta y + \gamma \sigma(x,y)\|_g.
\]

This norm measures the magnitude of the fused representation.

% ---------------------------------------------------------------------
\subsection*{5.4 Dominance Condition}

Nonlinear dominance occurs when
\[
\|\gamma \sigma(x,y)\|_g
\gg
\|\alpha x + \beta y\|_g.
\]

Equivalently,
\[
\frac{\|\gamma \sigma(x,y)\|_g}
     {\|\alpha x + \beta y\|_g}
\to \infty.
\]

Dominance indicates that nonlinear interaction overwhelms stable
blending, producing representational spikes.

% ---------------------------------------------------------------------
\subsection*{5.5 Fusion Instability Criterion}

Fusion instability occurs when
\[
\|F(x,y)\|_g > f_{\mathrm{thr}}.
\]

Spike instability occurs when
\[
\|\gamma \sigma(x,y)\|_g > f_{\mathrm{thr}}.
\]

Oscillatory instability occurs when
\[
F(x_{t+1}, y_{t+1}) \cdot F(x_t, y_t) < 0,
\]
in an appropriate operator sense.

Collapse instability occurs when
\[
F(x,y) \to 0
\]
due to destructive interference.

% ---------------------------------------------------------------------
\subsection*{5.6 Fusion–Field Dynamics Along Trajectories}

Given a latent trajectory $x(t)$ and a secondary trajectory $y(t)$,
fusion dynamics are defined by
\[
F_t = F(x(t), y(t)).
\]

Instability arises when $F_t$ exhibits:
\begin{itemize}
\item explosive growth,
\item collapse,
\item oscillation,
\item dominance of nonlinear interaction.
\end{itemize}

These behaviors correspond to distinct FFH subtypes.

% ---------------------------------------------------------------------
\subsection*{5.7 Fusion Stability Region}

Define the fusion stability region
\[
\Omega_f
= \{ (x,y) \in TM \times TM :
\|F(x,y)\|_g \le f_{\mathrm{thr}} \}.
\]

A latent trajectory pair $(x(t), y(t))$ is fusion–stable if
$(x(t), y(t)) \in \Omega_f$ for all $t$.

% ---------------------------------------------------------------------
\subsection*{5.8 Discrete Approximation}

In numerical settings, fusion–field output is computed directly from
discrete latent vectors:
\[
F_t = \alpha x_t + \beta y_t + \gamma \sigma(x_t, y_t).
\]

Nonlinear interaction terms may be implemented using:
\begin{itemize}
\item elementwise nonlinearities,
\item bilinear forms,
\item cross–product analogues,
\item neural activation functions.
\end{itemize}

The Python toolkit implements these approximations and validates them in
\texttt{/tests/test\_fusion.py}.

% ---------------------------------------------------------------------
\subsection*{5.9 Fusion–Field Hallucinations (FFH)}

A latent trajectory exhibits FFH if there exist sequences $x(t)$ and
$y(t)$ such that
\[
\|F(x(t), y(t))\|_g > f_{\mathrm{thr}}
\quad \text{for infinitely many } t.
\]

Equivalently, FFH occurs when nonlinear interaction prevents the latent
trajectory from remaining within a stable blending regime.

Fusion–field instability corresponds to uncontrolled nonlinear mixing of
representations.

% =====================================================================
% End of 05-Fusion.tex
% =====================================================================
